\section{Earned autonomy}

The 2026-07-28 release ships into a conformance-gated, SDK-tiered ecosystem: a
conformance suite live since 2026-01-23, published SDK tiers since 2026-02-23,
and, as of SEP-2484 (Final 2026-06-04), a standing rule that future
Standards-Track SEPs cannot reach Final without a matching conformance scenario
\citep{sep2484, sdktiers}. The SDK tiers are the part worth studying. A Tier 1
SDK must pass 100\% of the conformance suite; if any test fails continuously for
four weeks, it is relegated to Tier 2; and features move through a lifecycle with
a minimum twelve-month deprecation window measured from the revision release
\citep{sdktiers, mcplifecycle}. Standing is earned by demonstrated conformance
and lost publicly when the demonstration lapses.

That is a working model of earned standing, and it is a model humans run. No
automated controller relegates an SDK. Maintainers publish tiers against measured
conformance, on defined clocks, with a defined relegation rule, and the
governance works. This is the precise shape to port. Earned autonomy is a rolling
trust review in which workflow configurations hold published autonomy tiers,
delivery-evidence-layer signals are the input, promotion and relegation follow
defined triggers, and the resulting grant is enforced per call at the boundary
the spec provides. Steer is that cadence. It is not a controller that reads a
metric and rewrites policy in a closed loop; it is a review rhythm with a clock
and a rule, of the same kind the standard already runs on itself.

Every component of this cadence is mechanically expressible today. Identity and
per-call enforcement come from the spec. Attribution, outcome evaluation, and
served evidence come from the delivery evidence layer. Published tiers, clocks,
and triggers are process. The only genuinely new input is the swap at the center:
replace static scope grants, issued once and rarely revisited, with grants
conditioned on delivered-outcome evidence. Automating the cadence, turning the
rolling human review into a closed control loop, is the forward agenda, not a
claim about today.

Stated plainly, and stated as the sentence a skeptic should attack: autonomy
conditioned on task accuracy re-creates at the trust layer the same clean-record
trap the audit layer already has; delivered outcome is the only conditioning
variable that closes the loop, and the 2026-07-28 spec is the first standard
substrate on which that condition is enforceable per action.

\paragraph{The failure modes are dials, not objections.} Because Steer is a
cadence, its known failure modes are design parameters that set the cadence's
dials. \emph{Attribution lag} --- outcomes arrive days or weeks after the change
--- sets the length of the trailing window the review scores on; instantaneous
scoring is the mistake the window length corrects. \emph{Unstable identity} ---
agent instances are ephemeral and interchangeable --- sets the unit of trust: the
durable thing to govern is a workflow configuration, not an instance, because the
configuration is what produced the work and what will produce it again.
\emph{Goodhart's law} sets the gating signal: gate on held-in-production outcomes
and never on activity counts, because the moment the metric is a throughput count
it is gamed. \emph{Small-n} --- too few observations to trust --- sets the
cold-start default: absent evidence, a workflow starts in a sandbox tier, since
autonomy is earned from evidence and the safe default without it is low.

\paragraph{Two clarifications keep the framework honest} about what the protocol
does and does not provide. First, the per-call boundary is an enforcement point,
not a scope engine. The spec gives the point at which authorization is checked on
every call; it does not decide what a given principal may do. The scope
semantics, the policy itself, remain the implementer's design. Authentication is
not authorization scope. Second, per-call identity attributes to the
authenticated principal and the named client implementation, carried in
\texttt{clientInfo}, and not to an individual agent or subagent. Agent-level
attribution is a convention an organization builds on top, and it is exactly why
the unit of trust is a workflow configuration: that is the granularity the
evidence layer can attribute to, not a per-instance identity the protocol never
reported.

This shape has precedent outside software delivery. Open-source foundations
already promote projects by earned standing: the CNCF moves a project from
Sandbox to Incubating to Graduated on demonstrated adoption, security review, and
governance maturity \citep{cncflifecycle}. The criteria are partly quantitative
and partly qualitative, and the decision is ratified by humans. It is
evidence-based maturation with human ratification, not a purely metric-gated
switch, which is exactly the posture earned autonomy takes toward a workflow
rather than a whole project.

There is also direct empirical support for building the feedback capability the
framework depends on. DORA's 2025 research found that AI acts as an amplifier: it
has a positive relationship with delivery throughput and a negative one with
stability, and which effect dominates is conditional on a team's existing
capability \citep{dora2025}. The organizational property that decides whether AI
speed helps or hurts is the delivery-feedback capability, which is the same
capability this paper argues an organization must build. The amplifier result is
not only a stake. It is corroboration that the conditioning variable is the right
one.

The cadence is configuration, not research. The trust decision --- read the
delivered-outcome evidence, set the tier --- is the organization's to make; the
gateway's only job is to enforce the resulting scope on every call. A workflow
holds the identity its current tier has earned, and the gateway gates each tool
call against that identity:

\begin{figure}[t]
\centering
\begin{minipage}{0.92\linewidth}
\begin{lstlisting}[style=policy]
# Per-call authorization at an MCP gateway. A workflow holds the
# identity its current tier earns; the gateway enforces scope on
# every tool call. Relegation = the cadence reassigns the workflow
# to a lower-tier identity, and its next call is gated anew.
authorization:
  - allow: workflow == "wf-alpha"          # trusted: full access
  - allow: workflow == "wf-beta"           # sandbox: read-only
        && tool in ["read_file", "list_directory", "search_files"]
\end{lstlisting}
\end{minipage}
\caption{\textbf{Validated on a running MCP gateway.} A per-call authorization
rule gating tool scope on a workflow's tier-earned identity. On an isolated
gateway instance, the identical \texttt{write\_file} call was allowed for the
trusted identity and denied per call for the sandbox identity, while the read
call was allowed for both (see \texttt{data/}). The evidence-to-tier decision is
the organization's cadence; the gateway is only the actuator. The trailing window
the cadence scores on absorbs attribution lag, and the unit of trust is a workflow
configuration, not an instance. This is the shape teams already use to gate a
deploy, pointed at delivered-outcome evidence.}
\end{figure}
